% Notes of the Montreal seminar, 1970 (cote 129, typed part): archive cote 129, site page 71 (PDF page 72).
% typescript. Draft compared: none.
% AI-transcribed by Claude Opus 5.5 (claude-opus-5-5, adaptive thinking, effort high), Message Batches API, batch mtl129_full (msgbatch_01AGnDKHwy6ATCM7ehUa6263), 2026-11-13; not yet checked by eye.
\scanpage{129/71}
\mspage{\hand{III-1\unsure{6}}}

\[
X(S') = \varinjlim_{n} \mathrm{Hom}_S\ (S',\ \mathrm{Spec}\ \underline{\mathcal{O}}_S[[T_1,\ldots,T_r]]/(T_1,\ldots,T_r)^n)
\]
\[
= \varinjlim_{n} \{f_1^*,f_2,\ldots,f_r \in \Gamma(S',\underline{\mathcal{O}}_S)\ \mid\ f_i^n = 0,\ \ 1 \leq i \leq r\}
\]

2) $X$ est en particulier limite inductive de schémas finis
et radiciels au dessus de $S$.

On appelle \uline{groupe de Lie formel sur $S$} un faisceau en groupes
sur $S$ tel que sa section unité lui donne une structure de variété de
Lie formelle. Comme le produit de deux variétés de Lie formelles est
encore une variété de Lie formelle, il est équivalent de dire qu'un grou-%
pe de Lie formel est un objet en groupes de la catégorie des variétés de
Lie formelles sur $S$.

Dans ce qui suit, tous les groupes de Lie formels seront suppo-%
sés commutatifs.

Soit $G$ un groupe de Lie formel sur $S$. On vérifie aisément
que, si $p$ est localement nilpotent sur $S$, $G$ est de $p$-torsion, c'est-%
à-dire que
\[
G = G(\infty) = \varinjlim_{n} G(n) \ ,
\]
de sorte que $G$ est un groupe de Barsotti-Tate si et seulement si les
deux conditions suivantes sont vérifiées:

1) $G$ est $p$-divisible.

2) $G(1)$ est représentable par un schéma fini localement libre.

Il se trouve que \uline{si $S$ est artinien local}, d'unique point $s \in S$,

\mspage{\hand{71}}
